% Generated draft: requires mathematical and source review.
\documentclass[11pt]{article}
\usepackage[margin=25mm]{geometry}
\usepackage{fontspec}
\setmainfont{Latin Modern Roman}
\usepackage{amsmath,amssymb,mathtools,unicode-math}
\setmathfont{Latin Modern Math}
\usepackage{enumitem,needspace,xurl,hyperref}
\setlength{\parindent}{0pt}
\setlength{\parskip}{5pt}
\setlength{\emergencystretch}{3em}
\newcommand{\N}{\mathbb{N}}
\newcommand{\Z}{\mathbb{Z}}
\newcommand{\Q}{\mathbb{Q}}
\newcommand{\R}{\mathbb{R}}
\newcommand{\C}{\mathbb{C}}
\newcommand{\F}{\mathbb{F}}
\newcommand{\A}{\mathbb{A}}
\newcommand{\PP}{\mathbb{P}}
\newcommand{\E}{\mathbb{E}}
\newcommand{\eps}{\varepsilon}
\newcommand{\dd}{\,\mathrm{d}}
\begin{document}
\section*{UnsolvedMath ID: 701689}
\textbf{Source problem code:} catalogue entry
\subsection{Definitions and mathematical statement}
For a finite group \(G\), define its spectrum by
\[
\omega(G)=\{\operatorname{ord}(g):g\in G\},
\]
where \(\operatorname{ord}(g)\) is the least positive integer \(n\) such that \(g^n=1\). Define
\[
h(G)=\left|\left\{[H]_{\cong}: H\text{ is a finite group and }\omega(H)=\omega(G)\right\}\right|,
\]
where \([H]_{\cong}\) denotes the isomorphism class of \(H\). Thus groups are counted up to isomorphism; finiteness of this cardinality is not assumed.

Determine \(h(L)\) in each of the following cases:
\begin{enumerate}
\item[(a)] \(L=S_{10}\), the group of all permutations of a set with ten elements.
\item[(b)] \(L=\operatorname{Aut}(J_2)\), where \(J_2\) is the simple sporadic Janko group and \(\operatorname{Aut}(J_2)\) is its full automorphism group.
\item[(c)]
\[
\operatorname{PSL}(2,q)<L\leq\operatorname{Aut}(\operatorname{PSL}(2,q)).
\]
Here, for a prime power \(q\), \(\F_q\) denotes a field with \(q\) elements and
\[
\operatorname{PSL}(2,q)=\operatorname{SL}_2(\F_q)/Z(\operatorname{SL}_2(\F_q)),
\]
where \(\operatorname{SL}_2(\F_q)\) is the group of determinant-one matrices and \(Z\) denotes the center. The subgroup inclusion uses the identification of \(\operatorname{PSL}(2,q)\) with its group of inner automorphisms. The first inclusion is strict and the second permits equality. For every intended value of \(q\), determine \(h(L)\) for every such intermediate group \(L\).
\end{enumerate}
The supplied source does not specify the range of \(q\) in part (c). The finite-field notation requires a prime power, but whether the intended range includes \(q=2,3\), or is restricted to the simple cases \(q\geq4\), remains unresolved. This draft does not choose between those ranges.
\subsection{Short English statement}
Count, up to isomorphism, all finite groups having exactly the same element orders as \(S_{10}\), \(\operatorname{Aut}(J_2)\), and each group strictly above \(\operatorname{PSL}(2,q)\) inside its full automorphism group. The intended range of \(q\) requires clarification. The registry labels the problem open; that label and the cited status reference are supplied source information, not independently verified assessments.
\subsection{Sources}
\begin{itemize}
\item Local UnsolvedMath registry: permanent ID \(701689\), exact source problem code \texttt{catalogue entry}, title ``Kourovka Notebook Problem 20.31.'' The supplied record labels the problem \texttt{open} and records an update date of 2026-10-16.
\item \emph{The Kourovka Notebook (Unsolved Problems in Group Theory)}, problem 20.31, proposed by M. A. Grechkoseeva and A. V. Vasil'ev in 2022. The supplied background identifies the 20th issue (2022) and the 21st issue edition (2026), arXiv:1401.0300v48. The supplied collection description names E. I. Khukhro and V. D. Mazurov as editors.
\item The source directs readers to Subsection 4.2 of M. A. Grechkoseeva, V. D. Mazurov, W. J. Shi, A. V. Vasil'ev, and N. Yang, \emph{Commun. Math. Stat.} \textbf{11}, no. 2 (2023), 169--194, for the current status. No article title, URL, or contents of that subsection were supplied. This reference and its status qualification are reproduced from the record and have not been independently checked.
\end{itemize}
\end{document}
